\section{Évaluation des vulnérabilités de ResNet50}

Avant de déployer des mécanismes de défense sophistiqués, une question fondamentale se pose : à quel point notre modèle ResNet50, entraîné par transfer learning classique, est-il vulnérable aux attaques adversariales ? Cette évaluation empirique nous permettra d'établir une baseline de référence et de justifier l'investissement dans des techniques de défense avancées. Nous avons donc conçu un protocole d'évaluation complet, allant de la constitution du dataset jusqu'à la quantification précise des vulnérabilités.

\subsection{Constitution du dataset expérimental}

\subsubsection{Du dataset minimal à l'augmentation stratégique}

Notre point de départ était modeste : une centaine d'images collectées manuellement, réparties équitablement entre objets sûrs et dangereux. Si cette quantité peut suffire pour une preuve de concept, elle pose un problème majeur pour l'entraînement d'un réseau profond comme ResNet50 : le risque de surapprentissage est élevé. Pour résoudre cette limitation sans recourir à une collecte manuelle longue et coûteuse, nous avons conçu un système d'augmentation de données à quatre niveaux de complexité.

Notre classe \texttt{DangerousObjectDataset} implémente ce mécanisme d'augmentation en héritant du pattern PyTorch standard :

\begin{lstlisting}[language=Python, caption=Classe de chargement avec augmentation stratégique]
class DangerousObjectDataset(Dataset):
    def __init__(self, root_dir, split='train', transform=None):
        self.root_dir = root_dir
        self.transform = transform
        self.images = []
        self.labels = []

        # Parcours des dossiers safe/ et dangerous/
        for label, class_name in enumerate(['safe', 'dangerous']):
            class_dir = os.path.join(root_dir, split, class_name)
            for img_name in os.listdir(class_dir):
                self.images.append(os.path.join(class_dir, img_name))
                self.labels.append(label)

    def __getitem__(self, idx):
        image = Image.open(self.images[idx]).convert('RGB')
        if self.transform:
            image = self.transform(image)  # Augmentation à la volée
        return image, self.labels[idx]
\end{lstlisting}

L'approche adoptée calcule les transformations à la volée plutôt qu'en amont, ce qui économise l'espace disque tout en maximisant la variabilité des échantillons. Les quatre niveaux d'augmentation génèrent des variations contrôlées : 25\% d'images subissent des transformations légères (retournements, rotations modérées), 35\% des transformations moyennes (perspective, flou), 25\% des transformations fortes (rotations importantes), et 15\% des transformations spécifiquement conçues pour la résistance adversariale (postérisation, solarisation).

\subsubsection{Les transformations adversariales : au-delà de l'augmentation classique}

Cette dernière catégorie d'augmentation (15\% des images) mérite une attention particulière. Les transformations adversariales résistantes ne sont pas de simples augmentations esthétiques : elles modifient les images de manière à perturber les futures perturbations adversariales tout en préservant le contenu sémantique. La postérisation réduit le nombre de niveaux de couleur à 4 bits, atténuant ainsi les modifications pixel-niveau typiques des attaques. Le flou gaussien (noyau 5×5, $\sigma \in [1.0, 3.0]$) lisse les perturbations localisées, tandis que la solarisation et l'égalisation d'histogramme redistribuent les intensités de manière non linéaire.

Grâce à cette stratégie, notre dataset initial de 100 images a été multiplié par un facteur 15, atteignant 1536 images (split 1024/256/256). Cette augmentation permet d'entraîner ResNet50 sans surapprentissage visible, tout en introduisant une certaine robustesse naturelle via la diversité des transformations.

\subsection{Métriques d'évaluation de la robustesse}

Comment mesurer objectivement la vulnérabilité d'un modèle face aux attaques adversariales ? Nous avons adopté une approche à deux niveaux : les métriques classiques sur données propres d'une part, et les métriques de robustesse adversariale d'autre part.

\subsubsection{Performance sur données propres et robustesse adversariale}

L'évaluation sur données propres repose sur les métriques standards (accuracy, precision, recall, F1-score) calculées en mode évaluation PyTorch (\texttt{model.eval()} avec \texttt{torch.no\_grad()}), ce qui désactive le Dropout et fige la Batch Normalization pour une inférence déterministe.

Pour la robustesse adversariale, nous avons retenu trois métriques complémentaires. La \textit{robust accuracy} mesure le pourcentage d'exemples adversariaux correctement classés. La \textit{robustness degradation} capture la différence absolue entre l'accuracy sur données propres et celle sur données adversariales. Mais la métrique la plus révélatrice reste l'\textit{Attack Success Rate} (ASR), qui quantifie directement l'efficacité de l'attaque :

\begin{equation}
\text{ASR} = \frac{N_{\text{clean\_correct}} - N_{\text{adversarial\_correct}}}{N_{\text{clean\_correct}}}
\label{eq:asr_simple}
\end{equation}

Une ASR de 53\% signifie que plus de la moitié des prédictions correctes peuvent être inversées par des perturbations imperceptibles. C'est cette métrique qui guidera notre évaluation des vulnérabilités.

\subsection{Protocole d'attaque et évaluation baseline}

\subsubsection{Entraînement du modèle baseline}

Notre modèle baseline suit un protocole de transfer learning classique : ResNet50 pré-entraîné sur ImageNet, dont la couche de classification finale est remplacée pour s'adapter à notre tâche binaire. L'entraînement utilise Adam (learning rate 0.001) avec un scheduler réduisant le taux tous les 7 epochs, et un early stopping (patience 5) qui sauvegarde le meilleur modèle basé sur la validation accuracy.

Cette approche standard cache cependant une faiblesse critique : le modèle ne voit jamais d'exemples adversariaux durant l'apprentissage. Il optimise ses représentations uniquement pour des données propres, sans développer de défense contre les perturbations malveillantes. C'est précisément cette vulnérabilité que nos tests vont révéler.

\subsubsection{Implémentation des attaques adversariales}

Deux attaques de référence permettent d'évaluer la robustesse : FGSM (Fast Gradient Sign Method) et PGD (Projected Gradient Descent). FGSM, l'attaque single-step, calcule le gradient de la loss par rapport à l'image et ajoute une perturbation : $x_{adv} = x + \epsilon \cdot \text{sign}(\nabla_x L)$. Simple et rapide, elle sert de test minimal de robustesse.

PGD, en revanche, effectue plusieurs itérations de perturbations plus petites avec projection sur la boule $L_\infty$ à chaque étape. Cette approche itérative explore mieux l'espace des perturbations :

\begin{lstlisting}[language=Python, caption=Implémentation de l'attaque PGD]
def pgd_attack(model, images, labels, epsilon=0.031, alpha=0.008,
               num_iter=10):
    # Initialisation avec bruit aléatoire
    delta = torch.zeros_like(images).uniform_(-epsilon, epsilon)
    delta.requires_grad = True

    for _ in range(num_iter):
        outputs = model(images + delta)
        loss = F.cross_entropy(outputs, labels)
        loss.backward()

        # Gradient ascent + projection L_inf
        delta.data = (delta + alpha * delta.grad.sign()).clamp(-epsilon, epsilon)
        delta.grad.zero_()

    return (images + delta).clamp(0, 1)
\end{lstlisting}

Pour nos tests, nous avons configuré trois scénarios : FGSM ($\epsilon=8/255$), PGD standard ($\epsilon=8/255$, 10 itérations), et PGD fort ($\epsilon=16/255$, 20 itérations). Le paramètre $\alpha$ (taille du pas) est fixé à $\epsilon/4$.

\subsubsection{Pipeline d'évaluation complet}

Le pipeline d'évaluation orchestre systématiquement tous les tests : chargement du meilleur modèle baseline, puis exécution de quatre séries de tests (données propres, FGSM, PGD standard, PGD fort). Chaque test génère un dictionnaire de métriques permettant une analyse comparative directe et reproductible.

\subsection{Résultats de vulnérabilité du baseline}

\subsubsection{Une dichotomie révélatrice}

Les résultats révèlent un contraste saisissant. Sur données propres, le modèle baseline atteint 100\% d'accuracy sur les 15 échantillons de test, confirmant l'efficacité du transfer learning depuis ImageNet pour notre tâche binaire. Face à l'attaque FGSM ($\epsilon=8/255$), cette performance se maintient : l'Attack Success Rate reste à 0\%, suggérant que l'attaque single-step ne trouve pas les perturbations optimales.

Mais PGD révèle une vulnérabilité critique. Avec les mêmes $\epsilon=8/255$ mais via 10 itérations, l'accuracy chute à 46.7\% avec une ASR de 53.3\% : plus de la moitié des prédictions correctes sont inversées par des perturbations imperceptibles. L'attaque PGD forte ($\epsilon=16/255$, 20 itérations) aggrave la situation avec seulement 26.7\% d'accuracy résiduelle (ASR de 73.3\%).

\begin{table}[h]
\centering
\caption{Résultats de robustesse du modèle baseline}
\begin{tabular}{lccc}
\hline
\textbf{Configuration} & \textbf{Accuracy} & \textbf{ASR} & \textbf{Degradation} \\
\hline
Clean (baseline) & 100\% & - & - \\
FGSM ($\epsilon=8/255$) & 100\% & 0\% & 0\% \\
PGD ($\epsilon=8/255$, 10 iter) & 46.7\% & 53.3\% & 53.3\% \\
PGD ($\epsilon=16/255$, 20 iter) & 26.7\% & 73.3\% & 73.3\% \\
\hline
\end{tabular}
\label{tab:baseline_results}
\end{table}

\subsubsection{Interprétation : le paradoxe de la robustesse}

Cette vulnérabilité s'explique par le protocole d'entraînement. Le transfer learning standard optimise uniquement sur données propres, produisant des features efficaces mais fragiles face aux perturbations $L_\infty$. Le réseau n'a jamais été exposé à des exemples adversariaux, il n'a donc développé aucune défense.

Cette situation illustre un paradoxe fondamental : une performance parfaite sur données propres ne garantit aucune robustesse adversariale. Le modèle peut apprendre des features discriminantes mais non robustes, exploitant des corrélations statistiques qui s'effondrent sous perturbations. C'est précisément ce que révèle le contraste entre 100\% d'accuracy nominale et 46.7\% sous PGD.

\subsubsection{Certification de sécurité : verdict de non-production}

Pour évaluer la viabilité du déploiement, nous avons défini un système de certification basé sur trois critères : clean accuracy ≥ 95\%, FGSM robustness ≥ 80\%, et PGD robustness ≥ 80\%. Ces seuils reflètent les exigences d'un système déployé dans un environnement potentiellement hostile.

Le modèle baseline obtient un \textbf{Security Grade C} (niveau de menace MEDIUM). Il passe les critères de clean accuracy (100\%) et FGSM robustness (100\%), mais échoue dramatiquement sur PGD robustness (46.7\% << 80\%). Verdict : \textit{Production Ready: NO}.

Cette évaluation quantitative valide notre hypothèse initiale et justifie l'implémentation de mécanismes de défense avancés (Section 4.3).

\subsection*{Synthèse}

Ce chapitre a établi empiriquement la vulnérabilité critique du modèle ResNet50 baseline : bien qu'atteignant 100\% d'accuracy sur données propres, il s'effondre à 46.7\% sous PGD standard ($\epsilon=8/255$), avec une ASR de 53.3\%. Cette dégradation massive démontre que le transfer learning classique ne confère aucune robustesse adversariale intrinsèque.

L'analyse révèle également une hiérarchie claire : FGSM (single-step) échoue complètement, tandis que PGD (itératif) inverse plus de la moitié des prédictions. Les attaques itératives exploitent plus efficacement les vulnérabilités en explorant systématiquement l'espace des perturbations.

Le système de certification (grade C, statut "non-production") établit une référence claire que les modèles défendus devront surpasser. Les sections suivantes présenteront l'implémentation de défenses (Adversarial Training, Differential Privacy) visant à franchir les seuils de certification tout en maintenant une performance acceptable.

\begin{thebibliography}{9}

\bibitem{paszke2019pytorch}
Paszke, A., Gross, S., Massa, F., et al. (2019). \textit{PyTorch: An Imperative Style, High-Performance Deep Learning Library}. Advances in Neural Information Processing Systems, 32.

\bibitem{madry2018towards}
Madry, A., Makelov, A., Schmidt, L., Tsipras, D., \& Vladu, A. (2018). \textit{Towards Deep Learning Models Resistant to Adversarial Attacks}. International Conference on Learning Representations (ICLR).

\bibitem{goodfellow2014explaining}
Goodfellow, I. J., Shlens, J., \& Szegedy, C. (2014). \textit{Explaining and Harnessing Adversarial Examples}. arXiv preprint arXiv:1412.6572.

\end{thebibliography}
